\documentclass[12pt]{article}

\setlength{\parindent}{0pt}
\setcounter{secnumdepth}{3}
\usepackage[margin=1in]{geometry}

% packages
\usepackage{amsmath, amsfonts, amsthm, amssymb}
\usepackage{graphicx, float}
\usepackage{mathtools}
\usepackage{titlesec}
\usepackage{interval}
\usepackage{hyperref}
\usepackage{siunitx}
\usepackage{titling}
\usepackage{vwcol}
\usepackage{setspace}
\usepackage{empheq}
\usepackage{cancel}
\usepackage{esdiff}
\usepackage{multicol}
\usepackage{mdframed}
\usepackage{esdiff}
\usepackage{tikzsymbols}
\usepackage{multicol}
\usepackage{tikz}
\usepackage{varwidth}
\usepackage{pgfplots}
\usepackage{fancyhdr}
\usepackage{enumitem}
\usepackage{tcolorbox}

\intervalconfig {
	soft open fences
}

% header
\pagestyle{fancy}
\fancyhf{}
\lhead{Ethan Fan}
\rhead{AP Calculus BC}
\cfoot{\thepage}

% section format
\titleformat{\section}
  {\large \bfseries}
  {}
  {0em}
  {}
\titlespacing*{\section}{0pt}{0.5em}{0.3em}

\titleformat{\subsection}[runin]
  {\bfseries}
  {}
  {0em}
  {}
\titlespacing*{\subsection}{0pt}{0pt}{0em}

\titleformat{\subsubsection}[runin]
  {\itshape}
  {(\roman{subsubsection})}
  {0em}
  {}

% commands
\newcommand{\alignedintertext}[1]{%
  \noalign{%
    \vskip\belowdisplayshortskip
    \vtop{\hsize=\linewidth#1\par
    \expandafter}%
    \expandafter\prevdepth\the\prevdepth
  }%
}

\newcommand*{\paren}[1]{\ensuremath\left(#1\right)}
\newcommand*{\limit}[2][x]{\ensuremath{\displaystyle\lim_{#1\to#2}}}
\newcommand*{\Limit}[3][x]{\ensuremath{\displaystyle\lim_{#1\to#2}\left[#3\right]}}
\newcommand*{\deriv}[1][x]{\ensuremath{\dfrac{\mathrm{d}}{\mathrm{d}#1}}}
\newcommand*{\Deriv}[2][x]{\ensuremath{\dfrac{\mathrm{d}}{\mathrm{d}#1}\left[#2\right]}}
\newcommand{\dydx}{\dfrac{dy}{dx}}
\newcommand{\dydt}{\dfrac{dy}{dt}}
\newcommand{\dxdt}{\dfrac{dx}{dt}}
\newcommand*{\iinteg}[2][x]{\ensuremath{\displaystyle\int #2\;\mathrm{d}#1}}
\newcommand*{\dinteg}[4][x]{\ensuremath{\displaystyle\int_{#2}^{#3}#4\;\mathrm{d}#1}}
\newcommand*{\abs}[1]{\ensuremath{\left|#1\right|}}
\newcommand*{\eps}{\varepsilon}
\newcommand*{\real}{\ensuremath\mathbb{R}}
\newcommand*{\integer}{\ensuremath\mathbb{Z}}
\newcommand*{\posint}{\ensuremath\mathbb{Z^+}}
\newcommand*{\cbrt}[1]{\ensuremath{\sqrt[3]{#1}}}
\newcommand*{\lheqzero}{\ensuremath{\underset{\text{L'H}}{\overset{\left[\frac00\right]}{=}}}}
\newcommand*{\lheqinfty}{\ensuremath{\underset{\text{L'H}}{\overset{\left[\frac{\infty}{\infty}\right]}{=}}}}
\newcommand{\tor}{\text{ or }}
\newcommand{\floor}[1]{\lfloor #1 \rfloor}

\newcommand{\vem}{\vspace{0.5em}}
\newcommand{\example}[1]{\section{Example #1}}
\newcommand{\problem}[1]{\section{Problem #1}}
\newcommand{\subproblem}[1]{\subsection{(#1)} \,}

\DeclareMathOperator{\dne}{DNE}
\DeclareMathOperator{\sgn}{sgn}

\newcommand{\definition}[1]{\begin{tcolorbox}[colback=red!5!white,colframe=red!75!black,parbox=false] #1 \end{tcolorbox}}

\newcommand{\theorem}[2]{\begin{tcolorbox}[title={#1},colback=blue!5!white,colframe=blue!75!black,parbox=false] #2 \end{tcolorbox}}

\allowdisplaybreaks
% \postdisplaypenalty=100000

% document
\begin{document}

\vspace*{0cm}

% opening
\begin{center}
    {\Large \bfseries Problem Set 2} \\[0.5cm]
    {\large Approximating Accumulation} \\[0.35cm]
    {\normalsize September 7, 2026}
\end{center}

% problems

\problem{1}

\subproblem{a}
\[
\Delta t_1=5, \Delta t_2=10,\Delta t_3=5,\Delta t_4=15, \Delta t_5=10, \Delta t_6=15, \sum_{k=1}^{6}\Delta t_k=60
\]

\subproblem{b}
\begin{gather*}
    L=5\cdot 1.2+10\cdot 2+5\cdot 2.6+15\cdot 2.4+10\cdot 1.8 +15\cdot 2.2=126 L\\
    R=5\cdot 2+10\cdot 2.6+5\cdot 2.4+15\cdot 1.8+10\cdot 2.2 + 15\cdot 1.5= 119.5L\\
    \paren{\frac{L}{s}}(s)=L
\end{gather*}

\subproblem{c}
No, as the function is neither monotonically increasing nor decreasing. \vem

\subproblem{d}
\begin{gather*}
    T=\frac{1}{2}(126+119.5)=122.75L\\
    T=\frac{1}{2}(5(1.2+2)+10(2+2.6)+5 (2.6+2.4)+15(2.4+1.8)+10(1.8+2.2)+15(2.2+1.5))=122.75L
\end{gather*}
It assumes that the values increase/decrease linearly.\vem

\subproblem{e}
\begin{gather*}
    \Delta L=2.6L\\
    \Delta R=3.6L\\
    \Delta T=0.65L
\end{gather*}
The trapezoidal estimate performed the best.\vem

\subproblem{f}
We only have discrete data points here, rather than an actual algebraic function. We need actual functions to apply knowledge of antiderivatives. 

\problem{2}

\subproblem{a}
\begin{gather*}
    \Delta t=1, t_0=0,t_1=1,t_2=2,t_3=3,t_4=4\\
    P(t_0)=2,P(t_1)=9,P(t_2)=14,P(t_3)=17,P(t_4)=18\\
    P(c_k)\Delta t= (kW)(s)=(kJ)
\end{gather*}

\subproblem{b}
\begin{gather*}
    L_4=2+9+14+17=42 kJ\\
    R_4=9+14+17+18=58kJ
\end{gather*}

\subproblem{c}
\begin{gather*}
    P'(t)=8-2t\\
    t\in (0,4) \implies P'(t)>0
\end{gather*}
As $P(t)$ is monotonically increasing, $L_4$ is too small and $R_4$ is too large.
\[
L_4
<true \ energy<R_4\]

\subproblem{d}
\begin{gather*}
    M_4=5.75+11.75+15.75+17.75=51kJ\\
    T_4=\frac{1}{2}(42+58)=50 kJ
\end{gather*}

\subproblem{e}
\begin{gather*}
    \Delta L_4=\frac{26}{3}kJ\\
    \Delta R_4=\frac{22}{3}kJ\\
    \Delta M_4=\frac{1}{3}kJ\\
    \Delta T_4 =\frac{2}{3}kJ
\end{gather*}
$M_4$ and $T_4$ are significantly better, as they consider the average of both the left- and right-sided sums.\vem

\subproblem{f}
\[
P''(t)=-2
\]
The curve is concave up, and thus the midpoint always sits above the average, and hence $M_4>T_4$.\vem

\subproblem{g}
\begin{gather*}
    \frac{2M_4+T_4}{3}=\frac{2\cdot 51+50}{3}=\frac{152}{3}=exact
\end{gather*}
The weighted average is exactly the true value.

\problem{3}

\subproblem{a}
\[
i(t)=t^2-6t+5=(t-5)(t-1)
\]
The current reverses at $t=1,5.$ The battery is charging on $(0,1) \cup (5,6)$. The battery is discharging on $(1,5)$.\vem

\subproblem{b}
\[
t_0=0,t_1=1,t_2=2,t_3=3,t_4=4,t_5=5,t_6=6
\]

\subproblem{c}
\[
M_6=2.25+(-1.75)+(-3.75)+(-3.75)+(-1.75)+(2.25)=-6.5C
\]
The heights 2 to 5 have negative values.\vem

\subproblem{d}
The net change, as some heights are positive while others are negative.\vem

\subproblem{e}
\[
2(2.25+1.75+3.75)=15.5C
\]
The total charge is 15.5 C. The negative signs are removed, and in turn the magnitudes of the heights are added. If the reversal had fallen in the middle of a subinterval, the negative and positive components would partially cancel, leading to an inaccurate estimate.\vem

\subproblem{f}
\begin{gather*}
    E_{net}= \frac{1}{2}C\\
    E_{total}=\frac{1}{6}C
\end{gather*}

\subproblem{g}
The total charge $\dfrac{46}{3}$C should be logged. The net charge is useless, as it does not count towards the total battery usage. 

\problem{6}

\subproblem{a}
\begin{gather*}
    R_n=\boxed{\frac{4(2n^2+3n+1)}{n^2}+2}\\
    L_n=\boxed{\frac{4(2n^2-3n+1)}{n^2}+2}\\
    \lim_{x \to \infty} L_n=\lim_{x \to \infty} R_n =\boxed{10kJ}
\end{gather*}

\subproblem{b}
\begin{gather*}
    f(x)=x^2-x\\
    f\paren{\frac{3k}{n}}=\paren{\frac{3k}{n}+1}^2-\paren{\frac{3k}{n}+1}=\paren{\frac{3k}{n}}^2+\paren{\frac{3k}{n}}\\
    L_n=\frac{3}{n} \sum_{k=0}^{n-1} f\paren{\frac{3k}{n}+1}
    =\frac{3}{n} \sum_{k=0}^{n-1} \paren{\frac{3k}{n}}^2+\frac{3}{n}\sum_{k=0}^{n-1}\paren{\frac{3k}{n}}\\
    =\frac{27}{n^3} \sum_{k=0}^{n-1} k^2+\frac{9}{n^2}\sum_{k=0}^{n-1}k
    =\boxed{\frac{9(n-1)(3n-1)}{2n^2}}\\
    R_n=\frac{3}{n} \sum_{k=1}^{n} f\paren{\frac{3k}{n}+1}
    =\frac{3}{n} \sum_{k=1}^{n} \paren{\frac{3k}{n}}^2+\frac{3}{n}\sum_{k=1}^{n}\paren{\frac{3k}{n}}\\
    =\frac{27}{n^3} \sum_{k=1}^{n} k^2+\frac{9}{n^2}\sum_{k=1}^{n}k
    =\boxed{\frac{9(n+1)(3n+1)}{2n^2}}\\
    \lim_{x \to \infty}L_n=\boxed{\frac{27}{2}}\\
    \lim_{x \to \infty}R_n=\boxed{\frac{27}{2}}\\
\end{gather*}

\subproblem{c}
\begin{gather*}
    R_n=\frac{4(2n^2+3n+1)}{n^2}+2\\
    L_n=\frac{4(2n^2-3n+1)}{n^2}+2\\
    R_{10}=8+\frac{6}{5}+\frac{1}{25}+2=\boxed{11.24kJ}\\
    L_{10}=8-\frac{6}{5}+\frac{1}{25}+2=\boxed{8.84kJ}\\
    R_{100}=8+\frac{12}{100}+\frac{1}{2500}+2=\boxed{10.1204kJ}\\
    L_{100}=8-\frac{12}{100}+\frac{1}{2500}+2=\boxed{9.8804kJ}\\
\end{gather*}
Each error roughly decrease by a multiple of 10 when $n$ is multiplied by 10. The true value is trapped by $R_n$ and $L_n$ as the function is monotonically increasing. If the values are averaged, the error further decreases significantly.

\problem{7}

\subproblem{a}
\begin{gather*}
    f(x)=x^3\\
    R_n=\frac{b}{n} \sum_{k=1}^n f\paren{\frac{kb}{n}}=\frac{b^4}{n^4}\sum_{k=1}^n k^3=\frac{b^4}{n^4}\frac{n^2(n+1)^2}{4}=\boxed{\frac{b^4}{4}\paren{1+\frac{1}{n}}^2}\\
    \lim_{x \to \infty} R_n =\boxed{\frac{b^4}{4}}
\end{gather*}

\subproblem{b}
\begin{gather*}
    f(x)=x^2\\
    L_n=\frac{1}{n} \sum_{k=0}^{n-1}f\paren{\frac{k}{n}}=\frac{1}{n^3} \sum_{k=0}^{n-1} k^2=\boxed{\frac{(n-1)(2n-1)}{6n^2}}\\
    R_n=\frac{1}{n} \sum_{k=1}^{n}f\paren{\frac{k}{n}}=\frac{1}{n^3} \sum_{k=1}^{n} k^2=\boxed{\frac{(n+1)(2n+1)}{6n^2}}\\
    T_n=\boxed{\frac{2n^2+1}{6n^2}}\\
    f\paren{\frac{k}{n}+\frac{1}{2n}}=\paren{\frac{k}{n}+\frac{1}{2}}^2=\paren{\frac{k}{n}}^2+\paren{\frac{k}{n^2}}+\frac{1}{4n^2}\\
    M_n=\frac{1}{n}\sum_{k=0}^{n-1}f\paren{\frac{k}{n}+\frac{1}{2n}}=\frac{1}{n} \sum_{k=0}^{n-1} \paren{\paren{\frac{k}{n}}^2+\paren{\frac{k}{n^2}}+\frac{1}{4n^2}}\\
    =\frac{1}{n^3}\paren{\frac{(n-1)n(2n-1)}{6}+\frac{(n-1)n}{2}+\frac{n}{4}}\\
    =\frac{2(n-1)(2n-1)+6(n-1)+3}{12n^2}=\boxed{\frac{4n^2-1}{12n^2}}\\
\end{gather*}

\subproblem{c}
\begin{gather*}
    M_n=\frac{4n^2-1}{12n^2} \implies \frac{1}{3}-\frac{4n^2-1}{12n^2} = \frac{1}{12n^2}\\
    T_n=\frac{2n^2+1}{6n^2} \implies \frac{1}{3}-\frac{2n^2+1}{6n^2}=-\frac{1}{6n^2}
\end{gather*}
As $y=x^2$ is concave up, the slope is monotonically increasing. Thus, the secant line always sits above the curve, overestimating the area, while the midpoint line holds more area sitting below the curve, underestimating the area. \vem

\subproblem{d}
\[
\alpha =\frac{2}{3}, \beta =\frac{1}{3}\implies \frac{2}{3}M_n+\frac{1}{3}T_n=\frac{1}{3}
\]
The combination of $\alpha$ and $\beta$ also applies for Problem 2.\vem

\subproblem{e}
$R_n$ and $L_n$ decrease error for 50 percent, while $M_n$ and $T_n$ decrease error for 75 percent per doubling of $n$. I would therefore utilize the $M_n$ and $T_n$ rule as they are more accurate and efficient.










\end{document}

